% Options for packages loaded elsewhere
\PassOptionsToPackage{unicode}{hyperref}
\PassOptionsToPackage{hyphens}{url}
\PassOptionsToPackage{space}{xeCJK}
\documentclass[
  a4paper,
  ja=standard,
  xelatex]{bxjsarticle}
\usepackage{xcolor}
\usepackage{amsmath,amssymb}
\setcounter{secnumdepth}{-\maxdimen} % remove section numbering
\usepackage{iftex}
\ifPDFTeX
  \usepackage[T1]{fontenc}
  \usepackage[utf8]{inputenc}
  \usepackage{textcomp} % provide euro and other symbols
\else % if luatex or xetex
  \usepackage{unicode-math} % this also loads fontspec
  \defaultfontfeatures{Scale=MatchLowercase}
  \defaultfontfeatures[\rmfamily]{Ligatures=TeX,Scale=1}
\fi
\usepackage{lmodern}
\ifPDFTeX\else
  % xetex/luatex font selection
    \setmainfont[]{Noto Serif CJK JP}
    \setsansfont[]{Noto Sans CJK JP}
  \ifXeTeX
    \usepackage{xeCJK}
    \setCJKmainfont[]{Noto Serif CJK JP}
          \fi
  \ifLuaTeX
    \usepackage[]{luatexja-fontspec}
    \setmainjfont[]{Noto Serif CJK JP}
  \fi
\fi
% Use upquote if available, for straight quotes in verbatim environments
\IfFileExists{upquote.sty}{\usepackage{upquote}}{}
\IfFileExists{microtype.sty}{% use microtype if available
  \usepackage[]{microtype}
  \UseMicrotypeSet[protrusion]{basicmath} % disable protrusion for tt fonts
}{}
\makeatletter
\@ifundefined{KOMAClassName}{% if non-KOMA class
  \IfFileExists{parskip.sty}{%
    \usepackage{parskip}
  }{% else
    \setlength{\parindent}{0pt}
    \setlength{\parskip}{6pt plus 2pt minus 1pt}}
}{% if KOMA class
  \KOMAoptions{parskip=half}}
\makeatother
\usepackage{longtable,booktabs,array}
\usepackage{calc} % for calculating minipage widths
% Correct order of tables after \paragraph or \subparagraph
\usepackage{etoolbox}
\makeatletter
\patchcmd\longtable{\par}{\if@noskipsec\mbox{}\fi\par}{}{}
\makeatother
% Allow footnotes in longtable head/foot
\IfFileExists{footnotehyper.sty}{\usepackage{footnotehyper}}{\usepackage{footnote}}
\makesavenoteenv{longtable}
\setlength{\emergencystretch}{3em} % prevent overfull lines
\providecommand{\tightlist}{%
  \setlength{\itemsep}{0pt}\setlength{\parskip}{0pt}}
\usepackage{bookmark}
\IfFileExists{xurl.sty}{\usepackage{xurl}}{} % add URL line breaks if available
\urlstyle{same}
\hypersetup{
  hidelinks,
  pdfcreator={LaTeX via pandoc}}

\author{}
\date{}

\begin{document}

\textbf{宛先:}\\
日本銀行 総裁 殿\\
内閣総理大臣 高市早苗 殿

\textbf{件名:}\\
「緊急人道支援におけるステーブルコイン決済システムの導入に関する追加政策申請」\\
－（THP-CLEAR 技術提言：2025年10月14日提出版に対する補完・最終版）－

\textbf{本文:}\\
拝啓　平素より格別のご高配を賜り、厚く御礼申し上げます。

先般（令和七年十月十四日）に提出いたしました「THP-CLEAR（国家分散監査・決済基盤）導入提言書」に関し、\\
その後の検証結果および現地適用要件を踏まえた\textbf{最終補完資料}として、\\
添付「緊急人道支援におけるステーブルコイン決済システムの導入に関する政策申請書」を提出いたします。

本資料は、既存提言書に対し以下の要素を追加し、\\
国家インフラとしての安全性・即応性をより高い次元で保証するものです。

\begin{enumerate}
\def\labelenumi{\arabic{enumi}.}
\item
  Witnessログ分散の最終仕様（冗長率120--150％確定）
\item
  VoidCoin（価値ゼロ監査トークン）の正式定義と交換所制限条件
\item
  暗号方式（AES-256-GCM＋ISE）および冪等キーHの正式採用
\item
  フェーズ1～3の実証KPI定義および責任分界の明確化
\end{enumerate}

以上により、THP-CLEARは\textbf{国家級の人道・金融インフラとして即時運用可能な完成形}となります。\\
ご査収のうえ、適切な機関での審議・承認を賜れれば幸甚に存じます。

敬具

令和七年十月吉日

\textbf{The Horizon Protocol 編纂者}

\section{\texorpdfstring{\textbf{緊急人道支援におけるステーブルコイン決済システムの導入に関する政策申請書}}{緊急人道支援におけるステーブルコイン決済システムの導入に関する政策申請書}}\label{ux7dcaux6025ux4ebaux9053ux652fux63f4ux306bux304aux3051ux308bux30b9ux30c6ux30fcux30d6ux30ebux30b3ux30a4ux30f3ux6c7aux6e08ux30b7ux30b9ux30c6ux30e0ux306eux5c0eux5165ux306bux95a2ux3059ux308bux653fux7b56ux7533ux8acbux66f8}

\textbf{提出日:} 2025年10月21日 \textbf{宛先:} 関係省庁 御中

\subsection{\texorpdfstring{\textbf{1.
提言の趣旨と目的}}{1. 提言の趣旨と目的}}\label{ux63d0ux8a00ux306eux8da3ux65e8ux3068ux76eeux7684}

\subsubsection{\texorpdfstring{\textbf{1.1.
背景}}{1.1. 背景}}\label{ux80ccux666f}

近年の国際情勢の不安定化に伴い、紛争地域や大規模災害被災地への緊急人道支援活動において、支援金の迅速かつ透明性の高い送金が喫緊の課題となっています。従来の金融網は、インフラの破壊や政情不安により機能不全に陥るケースが多く、支援が末端まで届かない事態が頻発しています。

\subsubsection{\texorpdfstring{\textbf{1.2.
目的}}{1.2. 目的}}\label{ux76eeux7684}

本提言は、日本円に価値が裏付けられたステーブルコイン（JPYC等）を活用し、低コストかつ強靭な決済監査システム「THP-CLEAR」を導入することで、人道支援活動における決済の継続性と透明性を確保することを目的とします。本システムは、\textbf{平時・有事の双方で運用可能な軽量分散アーキテクチャ}として機能し、有事においても支援金の流れを正確に追跡可能にします。

\subsection{\texorpdfstring{\textbf{2.
システムの概要}}{2. システムの概要}}\label{ux30b7ux30b9ux30c6ux30e0ux306eux6982ux8981}

THP-CLEARは、\textbf{資金の保管・送金・為替取引を一切行わず、二者間の決済事実のみを不可逆的に記録する中立的監査レイヤ}です。そのため、資金決済法等の規制における「為替取引」には該当せず、シンプルな三層構造で迅速な導入が可能です。

\begin{longtable}[]{@{}
  >{\raggedright\arraybackslash}p{(\linewidth - 4\tabcolsep) * \real{0.1905}}
  >{\raggedright\arraybackslash}p{(\linewidth - 4\tabcolsep) * \real{0.2540}}
  >{\raggedright\arraybackslash}p{(\linewidth - 4\tabcolsep) * \real{0.5556}}@{}}
\toprule\noalign{}
\begin{minipage}[b]{\linewidth}\raggedright
層
\end{minipage} & \begin{minipage}[b]{\linewidth}\raggedright
役割
\end{minipage} & \begin{minipage}[b]{\linewidth}\raggedright
主な機能
\end{minipage} \\
\midrule\noalign{}
\endhead
\bottomrule\noalign{}
\endlastfoot
\textbf{MQ (受付層)} & 支援金決済リクエストの受付と整列 &
・管理機関と決済機関からの情報をペア化 ・取引IDと時刻を付与 \\
\textbf{WN (検証層)} & 取引内容の突合と検証 &
・二系統の情報の整合性を確認 ・電子署名により取引を確定 \\
\textbf{DB (保存層)} & 検証済み記録の永続保存 &
・監査証跡（Witnessログ）を記録 ・データの分散冗長化と自己修復 \\
\end{longtable}

電源は民生品のバッテリー付きACアダプタを想定しており、停電時にも一定時間の稼働が可能です。また、LGWAN（総合行政ネットワーク）への接続も考慮されています。

\subsection{\texorpdfstring{\textbf{3.
技術仕様ハイライト}}{3. 技術仕様ハイライト}}\label{ux6280ux8853ux4ed5ux69d8ux30cfux30a4ux30e9ux30a4ux30c8}

\subsubsection{\texorpdfstring{\textbf{3.1.
通信データ形式}}{3.1. 通信データ形式}}\label{ux901aux4fe1ux30c7ux30fcux30bfux5f62ux5f0f}

全トランザクションは、\textbf{AES-256-GCM外装（IV12バイト/TAG16バイト/CRC4バイト）で封緘した352バイト（2,816ビット）の完全固定長パケット}で処理されます。これにより、通信効率を最大化し、システムの処理負荷を極限まで低減します。

\subsubsection{\texorpdfstring{\textbf{3.2.
監査証跡（Witnessログ）と冪等性}}{3.2. 監査証跡（Witnessログ）と冪等性}}\label{ux76e3ux67fbux8a3cux8de1witnessux30edux30b0ux3068ux51aaux7b49ux6027}

検証が完了した取引は、\textbf{固定長128バイト}のレコードとして記録されます。全ての取引記録は、前掲の冪等キー\texttt{H}を主キー（Single
Source of Truth,
SSOT）としてデータベースに記録されるため、通信障害時にも二重記帳の心配がありません。

\subsubsection{\texorpdfstring{\textbf{3.3.
暗号資産等の取り扱い（VoidCoin）}}{3.3. 暗号資産等の取り扱い（VoidCoin）}}\label{ux6697ux53f7ux8cc7ux7523ux7b49ux306eux53d6ux308aux6271ux3044voidcoin}

本システムは、金銭価値の移転を伴わない物資支援の記録やシステムテストのため、「VoidCoin」という概念を導入します。これは\textbf{価値ゼロの「存在証明トークン」であり、投機的利用の余地はありません。}

\begin{itemize}
\item
  \textbf{取引制限:}
  \textbf{\texttt{ExchangeFlag=1}、事前登録された交換所ノードの関与、VoidCoin同士の交換禁止}という3つの厳格な条件下でのみ、ステーブルコインとの交換証跡を記録します。
\item
  \textbf{応用可能性:}
  将来的には、この仕組みを応用し、BTCやETHといった暗号資産による国際的な支援受け入れの監査にも対応可能です。
\end{itemize}

\subsection{\texorpdfstring{\textbf{4.
物理システム構成、コスト、運用設計}}{4. 物理システム構成、コスト、運用設計}}\label{ux7269ux7406ux30b7ux30b9ux30c6ux30e0ux69cbux6210ux30b3ux30b9ux30c8ux904bux7528ux8a2dux8a08}

（詳細は別紙B、Cをご参照ください）
汎用の省電力ミニPCを活用し、約2,000万円という低コストで100ノード規模のシステム導入が可能です。システムの技術的管理とネットワークインフラの管理責任は明確に分離されており、現地の負担を最小限に抑えます。

\subsection{\texorpdfstring{\textbf{5.
導入計画}}{5. 導入計画}}\label{ux5c0eux5165ux8a08ux753b}

（詳細は別紙Dをご参照ください）
リスクを管理しつつ着実な導入を図るため、「国内実証実験」「パイロット導入」「本格運用と拡大」の3段階のロードマップを提案します。

\subsection{\texorpdfstring{\textbf{6.
結論}}{6. 結論}}\label{ux7d50ux8ad6}

THP-CLEARは、最先端の技術を応用しつつも、「人道支援のための、低コストで運用が容易な決済監査システム」として設計されています。本システムの導入により、我が国が国際的な人道支援活動において主導的な役割を果たし、支援の透明性を世界に示す一助となることを確信しております。速やかなご検討と、実証実験への移行を提言いたします。

\section{\texorpdfstring{\textbf{別紙A：技術仕様詳細}}{別紙A：技術仕様詳細}}\label{ux5225ux7d19aux6280ux8853ux4ed5ux69d8ux8a73ux7d30}

\subsection{\texorpdfstring{\textbf{1.
設計思想}}{1. 設計思想}}\label{ux8a2dux8a08ux601dux60f3}

本システムの技術仕様は、「\textbf{軽量性・効率性・堅牢性}」を最優先としています。支援活動が現地の困難な状況下で行われることを想定し、最小の計算資源と低帯域ネットワークでも安定稼働するよう、全てのデータ形式をビットレベルで最適化しています。

\subsection{\texorpdfstring{\textbf{2.
通信パケット仕様（完全固定長：352バイト）}}{2. 通信パケット仕様（完全固定長：352バイト）}}\label{ux901aux4fe1ux30d1ux30b1ux30c3ux30c8ux4ed5ux69d8ux5b8cux5168ux56faux5b9aux9577352ux30d0ux30a4ux30c8}

全トランザクションは、\textbf{352バイト（2,816ビット）の単一パケット}で完結します。

\paragraph{\texorpdfstring{\textbf{◎
外装（暗号化層）：32バイト}}{◎ 外装（暗号化層）：32バイト}}\label{ux5916ux88c5ux6697ux53f7ux5316ux5c6432ux30d0ux30a4ux30c8}

通信の安全性を確保するため、\textbf{AES-256-GCM方式}で暗号化されます。

\begin{longtable}[]{@{}lll@{}}
\toprule\noalign{}
フィールド & 長さ (B) & 内容 \\
\midrule\noalign{}
\endhead
\bottomrule\noalign{}
\endlastfoot
\texttt{IV} & 12 & 暗号化 nonce \\
\texttt{TAG} & 16 & 認証タグ \\
\texttt{HdrCRC} & 4 & 破損検出用チェックサム \\
\end{longtable}

\paragraph{\texorpdfstring{\textbf{◎
本体（データ層）：320バイト}}{◎ 本体（データ層）：320バイト}}\label{ux672cux4f53ux30c7ux30fcux30bfux5c64320ux30d0ux30a4ux30c8}

取引の具体的な内容を格納します。

\begin{longtable}[]{@{}
  >{\raggedright\arraybackslash}p{(\linewidth - 6\tabcolsep) * \real{0.2500}}
  >{\raggedright\arraybackslash}p{(\linewidth - 6\tabcolsep) * \real{0.2500}}
  >{\raggedright\arraybackslash}p{(\linewidth - 6\tabcolsep) * \real{0.2500}}
  >{\raggedright\arraybackslash}p{(\linewidth - 6\tabcolsep) * \real{0.2500}}@{}}
\toprule\noalign{}
\begin{minipage}[b]{\linewidth}\raggedright
オフセット
\end{minipage} & \begin{minipage}[b]{\linewidth}\raggedright
長さ(B)
\end{minipage} & \begin{minipage}[b]{\linewidth}\raggedright
セクション
\end{minipage} & \begin{minipage}[b]{\linewidth}\raggedright
内容
\end{minipage} \\
\midrule\noalign{}
\endhead
\bottomrule\noalign{}
\endlastfoot
0 & 64 & \textbf{ヘッダ} &
取引ID、時刻、送金元/先ID、通貨種別、金額など \\
64 & 64 & \textbf{鍵ブロック} &
取引の正当性を証明するための暗号鍵（バイナリ形式） \\
128 & 128 & \textbf{署名ブロック} & 管理機関と決済機関双方の電子署名 \\
256 & 64 & \textbf{管理情報} &
\textbf{正式定義：冪等キー(\texttt{H})}、その他管理用データ \\
\end{longtable}

\subsection{\texorpdfstring{\textbf{3.
監査証跡仕様（Witnessログ：128バイト）}}{3. 監査証跡仕様（Witnessログ：128バイト）}}\label{ux76e3ux67fbux8a3cux8de1ux4ed5ux69d8witnessux30edux30b0128ux30d0ux30a4ux30c8}

検証が完了した取引は、\textbf{128バイトの固定長レコード}として記録されます。

\begin{longtable}[]{@{}lll@{}}
\toprule\noalign{}
オフセット & 長さ(B) & 内容 \\
\midrule\noalign{}
\endhead
\bottomrule\noalign{}
\endlastfoot
0 & 8 & 取引ID \\
8 & 8 & 取引時刻 \\
16 & 16 & 送金元/先ID \\
32 & 6 & 通貨種別コード \\
38 & 2 & ステータスフラグ \\
40 & 12 & 金額・手数料 \\
52 & 4 & 予約領域 \\
56 & 64 & 暗号鍵（再封緘済み） \\
120 & 8 & シーケンス番号 \\
\end{longtable}

\subsection{\texorpdfstring{\textbf{4.
制度運用を支える技術仕様}}{4. 制度運用を支える技術仕様}}\label{ux5236ux5ea6ux904bux7528ux3092ux652fux3048ux308bux6280ux8853ux4ed5ux69d8}

\paragraph{\texorpdfstring{\textbf{4.1.
冪等性の保証}}{4.1. 冪等性の保証}}\label{ux51aaux7b49ux6027ux306eux4fddux8a3c}

全ての取引には、以下の式で算出される\textbf{冪等キー\texttt{H}}が付与されます。
\textbf{\texttt{H\ =\ SHA3-256(TxID\ ‖\ Timestamp\_ISE)}}
データベースはこの\texttt{H}を主キーとして\texttt{UNIQUE}制約を課すことで、通信障害等で発生する重複トランザクションを自動的に吸収し、二重記帳を防止します。これがシステムにおける「単一の真実の源（Single
Source of Truth, SSOT）」となります。

\paragraph{\texorpdfstring{\textbf{4.2.
VoidCoinの検証ポリシー}}{4.2. VoidCoinの検証ポリシー}}\label{voidcoinux306eux691cux8a3cux30ddux30eaux30b7ux30fc}

VoidCoinに関する取引は、以下の全ての条件を満たす場合にのみシステムに受理され、制度的な安全性を担保します。

\begin{itemize}
\item
  \texttt{AssetCode}の最上位ビットが\texttt{1}であること。
\item
  \texttt{CcyNum}が\texttt{65535}、手数料(\texttt{FeeMinor})が\texttt{0}であること。
\item
  事前に登録された交換所ノードが関与し、\texttt{ExchangeFlag}が\texttt{1}の場合に限り、ステーブルコインとの交換証跡を記録する。
\item
  VoidCoin同士の直接交換はシステムレベルで拒否される。
\end{itemize}

\paragraph{\texorpdfstring{\textbf{4.3.
大口取引の内部処理}}{4.3. 大口取引の内部処理}}\label{ux5927ux53e3ux53d6ux5f15ux306eux5185ux90e8ux51e6ux7406}

取引金額は小数点以下18桁を64ビット整数(\texttt{i64})で扱いますが、理論上の上限を超える超大口取引（10¹⁸単位に近接）については、システム内部の計算時に自動的に128ビット整数(\texttt{u128})に拡張してオーバーフローを防ぎます。記録時は\texttt{i64}の範囲内に制約されます。

\section{\texorpdfstring{\textbf{別紙B：システム導入コスト試算}}{別紙B：システム導入コスト試算}}\label{ux5225ux7d19bux30b7ux30b9ux30c6ux30e0ux5c0eux5165ux30b3ux30b9ux30c8ux8a66ux7b97}

\subsection{\texorpdfstring{\textbf{1.
試算の前提}}{1. 試算の前提}}\label{ux8a66ux7b97ux306eux524dux63d0}

\begin{itemize}
\item
  本試算は、特定の紛争地域や被災地を対象とした\textbf{人道支援活動を想定した100ノード構成}の初期導入費用（CAPEX）および初年度運用費（OPEX）です。
\item
  ハードウェアは汎用品の活用を前提とし、調達コストを最小限に抑えます。
\item
  ソフトウェアはオープンソースを基盤としており、ライセンス費用は発生しません。
\end{itemize}

\subsection{\texorpdfstring{\textbf{2.
導入費用内訳（100ノード構成）}}{2. 導入費用内訳（100ノード構成）}}\label{ux5c0eux5165ux8cbbux7528ux5185ux8a33100ux30ceux30fcux30c9ux69cbux6210}

\begin{longtable}[]{@{}lllll@{}}
\toprule\noalign{}
項目 & 数量 & 単価（円） & 小計（円） & 備考 \\
\midrule\noalign{}
\endhead
\bottomrule\noalign{}
\endlastfoot
\textbf{サーバー機器 (CAPEX)} & & & \textbf{9,480,000} & HP Elite Mini
800 G9相当品 \\
MQ/WNノード (受付/検証) & 60 & 68,000 & 4,080,000 & 8GB RAM / 512GB
NVMe \\
DBノード (保存層) & 40 & 135,000 & 5,400,000 & 16GB RAM / 4TBx2
NVMe(RAID1) \\
\textbf{付属機材 (CAPEX)} & & & \textbf{600,000} & \\
バッテリー付きACアダプタ & 100 & 6,000 & 600,000 & 停電対策 \\
\textbf{通信・現地機材 (CAPEX)} & 一式 & & \textbf{7,000,000} & \\
可搬Wi-Fi中継装置 & 複数 & - & (3,000,000) & \\
衛星通信回線 (Starlink等) & 複数 & - & (2,000,000) & \\
現地用スマートフォン端末 & 200 & 10,000 & (2,000,000) & \\
\textbf{セットアップ・初年度運用費 (OPEX)} & 一式 & & \textbf{4,500,000}
& \\
保守・設置・技術者教育 & & - & (3,500,000) & \\
初年度通信・監査費用 & & - & (1,000,000) & \\
\textbf{小計（税抜）} & & & \textbf{21,580,000} & \\
\textbf{消費税（10\%）} & & & \textbf{2,158,000} & \\
\textbf{総計（税込）} & & & \textbf{23,738,000} & \\
\end{longtable}

\subsection{\texorpdfstring{\textbf{3.
コストに関する特記事項}}{3. コストに関する特記事項}}\label{ux30b3ux30b9ux30c8ux306bux95a2ux3059ux308bux7279ux8a18ux4e8bux9805}

\begin{itemize}
\item
  \textbf{価格変動リスク:}
  上記は2025年10月時点の国内市場価格に基づく参考見積もりです。実際の調達時には、\textbf{為替レート（±15\%程度）、国際物流費（±10\%程度）の変動リスク}を許容範囲として考慮する必要があります。
\item
  \textbf{価格交渉の余地:}
  本システムが人道支援における決済監査のデファクトスタンダードとなる可能性を背景に、\textbf{メーカー各社との交渉を通じて、さらなるコスト圧縮が期待できます。}
\item
  \textbf{圧倒的な低コスト:}
  約2,400万円という費用は、従来の金融システム構築・維持費用と比較して数桁低い水準であり、人道支援という予算が限られる領域において極めて高い費用対効果を発揮します。
\end{itemize}

\section{\texorpdfstring{\textbf{別紙C：運用設計と自己修復メカニズム}}{別紙C：運用設計と自己修復メカニズム}}\label{ux5225ux7d19cux904bux7528ux8a2dux8a08ux3068ux81eaux5df1ux4feeux5fa9ux30e1ux30abux30cbux30baux30e0}

\subsection{\texorpdfstring{\textbf{1.
責任分界}}{1. 責任分界}}\label{ux8cacux4efbux5206ux754c}

本システムの安定運用のため、システムそのものの管理責任と、通信インフラの管理責任を明確に分離します。これにより、専門性の高い領域にそれぞれが集中でき、人道支援の現場でも円滑な運用が可能となります。

\begin{longtable}[]{@{}
  >{\raggedright\arraybackslash}p{(\linewidth - 2\tabcolsep) * \real{0.5000}}
  >{\raggedright\arraybackslash}p{(\linewidth - 2\tabcolsep) * \real{0.5000}}@{}}
\toprule\noalign{}
\begin{minipage}[b]{\linewidth}\raggedright
管理主体
\end{minipage} & \begin{minipage}[b]{\linewidth}\raggedright
主な責任範囲
\end{minipage} \\
\midrule\noalign{}
\endhead
\bottomrule\noalign{}
\endlastfoot
\textbf{システム運用者} (THP-CLEARチーム) &
・全ノードの物理的な設置、ソフトウェアの設定・更新
・監査証跡（Witnessログ）の整合性管理
・データの分散冗長化と、障害発生時の自動復旧プロセスの維持
・暗号鍵の管理と監視端末の運用 \\
\textbf{ネットワーク管理者} (現地協力団体/自治体等) &
・各ノード間の通信経路（有線LAN、Wi-Fi、VPN等）の確保
・\textbf{衛星通信や無線といったバックホール回線の選定・維持管理}
・通信の安定性維持とセキュリティ設定 ・ノードの物理的な設置場所の提供 \\
\end{longtable}

\subsection{\texorpdfstring{\textbf{2.
データ保全と自己修復メカニズム}}{2. データ保全と自己修復メカニズム}}\label{ux30c7ux30fcux30bfux4fddux5168ux3068ux81eaux5df1ux4feeux5fa9ux30e1ux30abux30cbux30baux30e0}

本システムは、一部の機器が紛争や災害によって物理的に破壊・損失されることを前提として設計されており、以下の自律的な自己修復機能を備えています。

\begin{enumerate}
\def\labelenumi{\arabic{enumi}.}
\item
  \textbf{定期エクスポート:}
  各DBノードは、保持している監査証跡と鍵情報を\texttt{cron}（自動実行機能）により定期的にバックアップします。
\item
  \textbf{データの分散と暗号化:}
  バックアップデータは複数のかたまり（チャンク）に分割され、電子署名を付与された上で、他の複数のノードに\textbf{AES-256-GCM方式}で暗号化してコピーが配置されます。
\item
  \textbf{自動復元（再実体化）:}
  故障等で交換された新しいノードは、起動時に自身のTPM（セキュリティチップ）に記録された固有IDを元に、ネットワーク上から自身のデータのコピーを自動的に収集し、正常な状態に復元します。
\end{enumerate}

この仕組みにより、\textbf{人為的な操作をほぼ介さずにシステムの完全性が維持}され、過酷な環境下での運用を可能にします。また、広域クラウド障害の教訓を踏まえ、\textbf{複数クラウド事業者と自治体内ネットワークに拠点を分散させることで、通信断絶時も継続稼働できる}設計となっています。

\subsection{\texorpdfstring{\textbf{3.
監視端末と鍵データベースの運用}}{3. 監視端末と鍵データベースの運用}}\label{ux76e3ux8996ux7aefux672bux3068ux9375ux30c7ux30fcux30bfux30d9ux30fcux30b9ux306eux904bux7528}

\begin{itemize}
\item
  \textbf{役割:}
  システム全体の稼働状況を監視し、全取引の正当性を担保する「鍵データベース」を安全に保管します。
\item
  \textbf{配置の柔軟性:}
  システムと同一ネットワーク上にあれば、物理的な場所を問わず運用可能です。遠隔地からの安全な監視を実現します。
\item
  \textbf{鍵データベースの安全性:}

  \begin{itemize}
  \item
    取引毎の暗号鍵は、さらに別の暗号鍵で何重にも保護された状態で保管されます。
  \item
    \textbf{個人情報非保持:}
    鍵データベースは個人特定情報（PII）を保持せず、発行主体ドメインと内部IDのみで管理されるため、プライバシーリスクを最小化します。
  \item
    \textbf{単一障害点（SPOF）の排除:}
    鍵データベース自体も自己修復メカニズムの対象となっており、監視端末が万一故障しても、他のノードから自動的に復元されます。
  \end{itemize}
\end{itemize}

これらの運用設計により、THP-CLEARは最小限の人的リソースで、国家レベルのセキュリティと信頼性を人道支援の現場へ提供します。

\section{\texorpdfstring{\textbf{別紙D：リスク評価、ガバナンス、及び導入計画}}{別紙D：リスク評価、ガバナンス、及び導入計画}}\label{ux5225ux7d19dux30eaux30b9ux30afux8a55ux4fa1ux30acux30d0ux30caux30f3ux30b9ux53caux3073ux5c0eux5165ux8a08ux753b}

\subsection{\texorpdfstring{\textbf{1.
リスク評価と対策}}{1. リスク評価と対策}}\label{ux30eaux30b9ux30afux8a55ux4fa1ux3068ux5bfeux7b56}

本システムは、\textbf{平時・有事の双方で運用可能な軽量分散アーキテクチャ}として機能します。その設計思想である「シンプルさと堅牢性」により、多くのリスクは構造的に低減されています。

\begin{longtable}[]{@{}
  >{\raggedright\arraybackslash}p{(\linewidth - 4\tabcolsep) * \real{0.0486}}
  >{\raggedright\arraybackslash}p{(\linewidth - 4\tabcolsep) * \real{0.1597}}
  >{\raggedright\arraybackslash}p{(\linewidth - 4\tabcolsep) * \real{0.7917}}@{}}
\toprule\noalign{}
\begin{minipage}[b]{\linewidth}\raggedright
リスク分類
\end{minipage} & \begin{minipage}[b]{\linewidth}\raggedright
想定される事象
\end{minipage} & \begin{minipage}[b]{\linewidth}\raggedright
主な対策
\end{minipage} \\
\midrule\noalign{}
\endhead
\bottomrule\noalign{}
\endlastfoot
\textbf{技術的リスク} & ・ハードウェアの故障 ・ソフトウェアの不具合
・現地の通信網の不安定化 ・大規模クラウド障害 &
・\textbf{機器の冗長化と自己修復機能:}
一部のノードが故障しても、システム全体は停止せず、データは自動的に復元されます。
・\textbf{シンプルな構成:}
汎用部品の使用により、現地での代替機交換が容易です。
・\textbf{低帯域設計:}
衛星通信のような不安定な回線でも動作するよう、通信データ量を極小化しています。
・\textbf{クラウド非依存:}
広域クラウド障害の教訓を活かし、複数クラウド事業者と自治体内ネットワークに拠点を分散させ、単一障害点を排除します。 \\
\textbf{運用上のリスク} & ・暗号鍵の不適切な管理 ・現地担当者の操作ミス
・ノードの物理的な盗難・紛失 & ・\textbf{責任分界の明確化:}
システムの技術的管理は専門チームが遠隔で実施し、現地担当者の作業を最小限にします。
・\textbf{自律運用:}
ノードは電源を入れるだけで自動的に機能するため、専門的な操作は不要です。
・\textbf{データの暗号化:}
ノード内のデータは全て高度に暗号化されており、万一盗難されても情報が漏洩することはありません。
・\textbf{個人情報非保持:}
鍵データベースは個人特定情報（PII）を保持せず、発行主体ドメインと内部IDのみで管理します。 \\
\textbf{セキュリティリスク} & ・サイバー攻撃による通信妨害
・不正な取引記録の試み & ・\textbf{独立した通信網:}
公共インターネットから隔離された閉域網（VPN等）での運用を基本とします。
・\textbf{強力な暗号化と署名:}
通信は\textbf{AES-256-GCM方式}で保護され、全ての取引は複数の電子署名によって検証されるため、不正な記録はシステムに受理されません。 \\
\textbf{法的・地政学的リスク} & ・支援対象国の法規制との抵触
・政情不安による活動の制約 & ・\textbf{中立的な監査機能:}
本システムは、資金の保管・送金・為替取引を一切行わず、二者間の決済事実のみを不可逆的に記録する\textbf{中立的監査レイヤ}です。そのため、資金決済法等の規制における「為替取引」には該当しません。
・\textbf{国際NGOとの連携:}
現地の法制度や情勢に詳しい国際的な人道支援団体と連携し、コンプライアンスを確保します。 \\
\end{longtable}

\subsection{\texorpdfstring{\textbf{2.
ガバナンス・運用体制案}}{2. ガバナンス・運用体制案}}\label{ux30acux30d0ux30caux30f3ux30b9ux904bux7528ux4f53ux5236ux6848}

本システムの公共性と透明性を担保するため、以下の多層的なガバナンス・運用体制を提案します。

\begin{longtable}[]{@{}
  >{\raggedright\arraybackslash}p{(\linewidth - 4\tabcolsep) * \real{0.3333}}
  >{\raggedright\arraybackslash}p{(\linewidth - 4\tabcolsep) * \real{0.3333}}
  >{\raggedright\arraybackslash}p{(\linewidth - 4\tabcolsep) * \real{0.3333}}@{}}
\toprule\noalign{}
\begin{minipage}[b]{\linewidth}\raggedright
階層
\end{minipage} & \begin{minipage}[b]{\linewidth}\raggedright
主体
\end{minipage} & \begin{minipage}[b]{\linewidth}\raggedright
主な役割
\end{minipage} \\
\midrule\noalign{}
\endhead
\bottomrule\noalign{}
\endlastfoot
\textbf{監督・意思決定層} &
\textbf{関係省庁および国際機関による合同評議会} &
・システム全体の運用方針の決定 ・支援対象とする活動の承認
・予算の監督と活動評価 \\
\textbf{技術・運用層} & \textbf{システム運用者 (THP-CLEARチーム)} &
・評議会の決定に基づくシステムの技術的な維持管理 ・24時間365日の稼働監視
・障害発生時の復旧対応と技術サポート \\
\textbf{現地協力層} & \textbf{国際NGO・現地協力団体} &
・ネットワーク管理者として、現地での通信インフラの確保（衛星・無線バックホールの選定を含む）
・ノードの物理的な設置場所の提供と管理 ・現地での支援活動の実行 \\
\end{longtable}

\subsection{\texorpdfstring{\textbf{3.
段階的な導入計画（ロードマップ）}}{3. 段階的な導入計画（ロードマップ）}}\label{ux6bb5ux968eux7684ux306aux5c0eux5165ux8a08ux753bux30edux30fcux30c9ux30deux30c3ux30d7}

本システムを安全かつ着実に導入するため、以下の3段階のロードマップを提案します。

\paragraph{\texorpdfstring{\textbf{フェーズ1：国内実証実験
(T+0～3ヶ月)}}{フェーズ1：国内実証実験 (T+0～3ヶ月)}}\label{ux30d5ux30a7ux30fcux30ba1ux56fdux5185ux5b9fux8a3cux5b9fux9a13-t03ux30f6ux6708}

\begin{itemize}
\item
  \textbf{目的:}
  システムの基本機能と安全性を国内の管理された環境で最終検証する。
\item
  \textbf{主要成功基準 (Key Performance Indicators):}

  \begin{itemize}
  \item
    \textbf{p99 ACK遅延:} 300秒未満
  \item
    \textbf{重複取引率:} 0.1\%未満（管理下での障害試験時）
  \end{itemize}
\end{itemize}

\paragraph{\texorpdfstring{\textbf{フェーズ2：パイロット導入
(T+3～6ヶ月)}}{フェーズ2：パイロット導入 (T+3～6ヶ月)}}\label{ux30d5ux30a7ux30fcux30ba2ux30d1ux30a4ux30edux30c3ux30c8ux5c0eux5165-t36ux30f6ux6708}

\begin{itemize}
\item
  \textbf{目的:}
  特定の地域・目的を定めた実際の人道支援活動に試験的に導入し、実環境での有効性を評価する。
\item
  \textbf{主要成功基準 (Key Performance Indicators):}

  \begin{itemize}
  \item
    \textbf{p99 ACK遅延:} 180秒未満
  \item
    \textbf{重複取引率:} 0.5\%未満（平常時）
  \item
    \textbf{データ復旧時間:} 10GBあたり25分未満
  \item
    \textbf{日次監査不整合件数:} 0件
  \item
    \textbf{大規模通信断（5時間）からの遅延解消:} 24時間以内
  \end{itemize}
\end{itemize}

\paragraph{\texorpdfstring{\textbf{フェーズ3：本格運用と拡大
(T+6ヶ月以降)}}{フェーズ3：本格運用と拡大 (T+6ヶ月以降)}}\label{ux30d5ux30a7ux30fcux30ba3ux672cux683cux904bux7528ux3068ux62e1ux5927-t6ux30f6ux6708ux4ee5ux964d}

\begin{itemize}
\item
  \textbf{目的:}
  パイロット導入の結果に基づき、システムを本格稼働させ、他の人道支援活動への適用を拡大する。
\item
  \textbf{内容:}

  \begin{enumerate}
  \def\labelenumi{\arabic{enumi}.}
  \item
    本システムを正式な人道支援インフラとして位置づけ、継続的な運用体制を確立。
  \item
    他の国際NGOや支援団体からの利用申請を受け付け、適用範囲を広げる。
  \end{enumerate}
\end{itemize}

\end{document}
